%% GERADO por gera_tab_compostos.py -- nao editar a mao
\begin{table}[!ht]
\centering
\caption{The 19 attribute combinations the summarization literature publishes, with the number of archive papers that use each. These are not ensembles assembled here: each row is a set of terms some published work combines. \emph{Sens.} is the within-reference mean against the degradation level on KADID-10k under the $z$-score aggregation. \textbf{Orientation here comes from the data}, the sign of the mean across types, because a combination of attributes has no direction declared by any source, whereas Table~S18 orients each operation by the direction its own source states. Magnitudes match across the two tables; a single-attribute row therefore carries the opposite sign to Table~S18 exactly when the declared direction contradicts the data. \textbf{This makes the gate here the more permissive of the two}, since $|\rho|\ge\rho^{*}$, so the verdict below is conservative: no composite clears it even under the looser orientation. \emph{Types} is the count above $.50$ of 25; \emph{Gen.} is the direction-aligned correlation on KonIQ-10k as a magnitude, under the same orientation rule as the sensitivity column. \textbf{Sensitivity and Gen.\,($z$) are the same configuration}; the ridge column is the generous variant, fitted outside the reference and reported separately rather than mixed into the verdict. \textbf{Eligibility below uses the $z$-score configuration throughout}, and the count is the same under either consistent choice: zero of the 19 are eligible whether both gates use the $z$ form or both use the ridge form. \textbf{1 of the 19 clears the sensitivity gate} (green) and none is eligible, because that one does not transfer.}
\label{tab:compostos}
\footnotesize
\setlength{\tabcolsep}{3pt}
\begin{tabular}{@{}>{\raggedright\arraybackslash}p{4.9cm}cccccc@{}}
\toprule
\rowcolor{headerblue}
\hd{Combination} & \hd{Papers} & \hd{Sens.} & \hd{Types} & \hd{Gen. ($z$)} & \hd{Gen. (ridge)} & \hd{Eligible} \\
\midrule
colourfulness + Laplacian variance + LBP uniformity & 1 & \textcolor{accentgreen}{\textbf{.524}} & 17 & .061 & .067 & no \\
Laplacian variance + LBP uniformity & 1 & .439 & 15 & .080 & .082 & no \\
edge distribution & 2 & .298 & 9 & .099 & .099 & no \\
spatial information & 5 & .251 & 13 & .513 & .513 & no \\
LBP uniformity + gradient magnitude & 1 & .189 & 8 & .177 & .007 & no \\
colourfulness & 2 & .182 & 9 & .007 & .007 & no \\
LBP uniformity & 4 & .175 & 12 & .205 & .205 & no \\
colourfulness + exposure & 1 & .173 & 10 & .243 & .003 & no \\
edge distribution + hue count + Laplacian variance & 1 & .156 & 7 & .204 & .347 & no \\
colourfulness + edge distribution + hue count + saliency & 2 & .148 & 7 & .162 & .117 & no \\
exposure + Laplacian variance & 1 & .108 & 9 & .500 & .357 & no \\
colourfulness + gradient magnitude & 1 & .108 & 11 & .296 & .066 & no \\
exposure & 1 & .102 & 4 & .343 & .343 & no \\
GLCM contrast + entropy & 1 & .092 & 11 & .371 & .084 & no \\
LBP uniformity + saliency & 3 & .089 & 9 & .071 & .277 & no \\
saliency & 6 & .075 & 3 & .333 & .333 & no \\
Laplacian variance & 3 & .047 & 9 & .601 & .601 & no \\
colourfulness + edge distribution & 1 & .017 & 4 & .085 & .051 & no \\
LBP uniformity + entropy & 1 & .009 & 8 & .007 & .285 & no \\
\bottomrule
\end{tabular}
\end{table}
